%=====================================================================
\chapter{Adversarial Search and Game Playing}\label{ch:adversarial}
%=====================================================================
\locator{3}{Part III --- Constraints, Games and Plans}

\begin{outcomes}
By the end of this chapter you will be able to:
\begin{tight}
  \item \textbf{Model} a two-party interaction as a game tree with a utility at
        each leaf, and \textbf{say} what the zero-sum assumption commits you to.
  \item \textbf{Compute} the minimax value of a position by hand and in code.
  \item \textbf{Apply} alpha--beta pruning, \textbf{trace} its cutoffs, and
        \textbf{state} the bound it achieves under perfect move ordering.
  \item \textbf{Explain} why move ordering determines how much is pruned, and
        \textbf{judge} when the cost of ordering exceeds what it saves.
  \item \textbf{Cut off} search with an evaluation function, and \textbf{recognise}
        the horizon effect.
\end{tight}
\end{outcomes}

\begin{prereq}
\chref{uninformed} (depth-first search, which is what a game tree search is) and
\chref{agents} --- specifically the \emph{multi-agent} property of a task
environment. This chapter is what you do when the sixth property of
Table~\ref{tab:envprops} comes out multi-agent.
\end{prereq}

\begin{keyterms}
game tree $\bullet$ ply $\bullet$ terminal state $\bullet$ utility $\bullet$
zero-sum $\bullet$ minimax value $\bullet$ MAX and MIN $\bullet$ alpha--beta
pruning $\bullet$ cutoff $\bullet$ move ordering $\bullet$ principal variation
$\bullet$ cutoff test $\bullet$ evaluation function $\bullet$ horizon effect
\end{keyterms}

\section{When the Other Party Is Choosing Too}

The delivery team and the client are agreeing a release slot. We propose, they
respond, we respond, and so on until somebody accepts or the deadline forces a
decision. Every option we might take is worth something to us --- and what we
eventually get depends on what they do next, which depends on what we do now.

\chref{local} would treat their behaviour as part of the landscape. \chref{csp}
would treat it as a constraint. Both are wrong in the same way: the other party is
not weather and not a rule. They are optimising, and if we model them as noise they
will exploit us systematically. That is the \emph{multi-agent} property of
\chref{agents}, and this chapter is what it costs to take it seriously.

\begin{definitionbox}[Game tree and minimax value]
A \term{game tree} has the current position at its root; the children of a node are
the positions reachable by one legal move; and each \term{terminal} position has a
\term{utility} for us. Levels alternate between \term{MAX} (our move --- we choose
the largest value) and \term{MIN} (their move --- they choose the smallest). One
level is a \term{ply}.

The \term{minimax value} of a node is what the position is worth under optimal play
by both sides:
\[
  \mathrm{minimax}(n) =
  \begin{cases}
    \mathrm{utility}(n) & \text{$n$ terminal} \\
    \max_{s} \mathrm{minimax}(s) & \text{$n$ is a MAX node} \\
    \min_{s} \mathrm{minimax}(s) & \text{$n$ is a MIN node}
  \end{cases}
\]
\end{definitionbox}

\begin{alertbox}[What MIN assumes, and when it is wrong]
Minimax assumes the opponent plays the move that is worst for us --- that the game
is \term{zero-sum}, our gain being exactly their loss. A negotiation usually is not:
both sides can lose by failing to agree, and both can gain from a slot that suits
each.

Assuming an optimal adversary is nevertheless the right default, because it is the
\emph{safe} assumption. The minimax value is a guarantee: we can obtain at least
this much whatever they do. Model them as cooperative and you get a better number
and no guarantee at all. Know which one you are quoting.
\end{alertbox}

\section{Minimax by Hand}

\dsref{game-tree} is a three-ply negotiation, small enough to compute in your head
and used for the rest of this chapter. Leaves are days of delay avoided, from our
point of view.

\dsfig{%
\begin{tikzpicture}[font=\scriptsize,
  mx/.style={rectangle, rounded corners=2pt, draw=primary, line width=1pt,
             fill=primary!12, text=primary!60!black, font=\scriptsize\bfseries,
             minimum width=0.72cm, minimum height=0.56cm},
  mn/.style={circle, draw=accent, line width=1pt, fill=accent!12,
             text=accent!60!black, font=\scriptsize\bfseries,
             minimum size=0.66cm, inner sep=0pt},
  lf/.style={rectangle, rounded corners=2pt, draw=neutral, line width=0.8pt,
             fill=gray!8, text=neutral!60!black, font=\scriptsize,
             minimum width=0.62cm, minimum height=0.52cm},
  ed/.style={neutral!70, line width=0.8pt},
  lab/.style={font=\tiny, text=neutral}]

\node[mx] (r) at (0,3.1) {5};
\node[mn] (m1) at (-2.6,1.9) {5};
\node[mn] (m2) at (2.6,1.9)  {0};
\node[mx] (a) at (-3.9,0.7) {5};
\node[mx] (b) at (-1.3,0.7) {9};
\node[mx] (c) at (1.3,0.7)  {2};
\node[mx] (d) at (3.9,0.7)  {0};
\node[lf] (l1) at (-4.55,-0.5) {3};
\node[lf] (l2) at (-3.25,-0.5) {5};
\node[lf] (l3) at (-1.95,-0.5) {6};
\node[lf] (l4) at (-0.65,-0.5) {9};
\node[lf] (l5) at (0.65,-0.5)  {1};
\node[lf] (l6) at (1.95,-0.5)  {2};
\node[lf] (l7) at (3.25,-0.5)  {0};
\node[lf] (l8) at (4.55,-0.5)  {-1};

\foreach \a/\b in {r/m1,r/m2,m1/a,m1/b,m2/c,m2/d,
                   a/l1,a/l2,b/l3,b/l4,c/l5,c/l6,d/l7,d/l8}{
  \draw[ed] (\a) -- (\b);}

\node[lab, anchor=east] at (-5.1,3.1)  {\textbf{MAX} (us)};
\node[lab, anchor=east] at (-5.1,1.9)  {\textbf{MIN} (them)};
\node[lab, anchor=east] at (-5.1,0.7)  {\textbf{MAX} (us)};
\node[lab, anchor=east] at (-5.1,-0.5) {utility to us};

\node[font=\tiny\itshape, text=neutral, text width=4.2cm, align=flush left,
      anchor=west] at (5.2,1.4)
  {Working upward: the MAX row takes $\max(3,5)=5$, $\max(6,9)=9$,
   $\max(1,2)=2$, $\max(0,-1)=0$. The MIN row takes $\min(5,9)=5$ and
   $\min(2,0)=0$. The root takes $\max(5,0)=5$.};
\node[font=\tiny\bfseries, text=primary, text width=4.2cm, align=flush left,
      anchor=west] at (5.2,-0.2)
  {So the position is worth 5, and our first move is to the left.};
\end{tikzpicture}}%
{A three-ply negotiation. Squares are our choices, circles theirs. Minimax
evaluates all eight leaves to establish that the position is worth $5$.}{game-tree}

\section{Alpha--Beta Pruning}

Minimax examined every leaf. Most of them could not have changed the answer.

\begin{definitionbox}[Alpha and beta]
Carry two bounds down the tree:
\begin{tight}
  \item $\alpha$ --- the best value MAX can already guarantee somewhere above.
  \item $\beta$ --- the best value MIN can already guarantee somewhere above.
\end{tight}
At any node, if $\alpha \ge \beta$, stop examining the remaining children: the
player above already has a better alternative and will never choose this branch, so
its exact value cannot matter.

\term{Alpha--beta pruning} returns exactly the minimax value. It is not an
approximation --- it computes the same answer while declining to evaluate leaves
that provably cannot affect it.
\end{definitionbox}

\dsref{alphabeta} traces it on the same tree. Five of eight leaves are evaluated.

\dsfig{%
\begin{tikzpicture}[font=\scriptsize,
  mx/.style={rectangle, rounded corners=2pt, draw=primary, line width=1pt,
             fill=primary!12, text=primary!60!black, font=\scriptsize\bfseries,
             minimum width=0.72cm, minimum height=0.56cm},
  mn/.style={circle, draw=accent, line width=1pt, fill=accent!12,
             text=accent!60!black, font=\scriptsize\bfseries,
             minimum size=0.66cm, inner sep=0pt},
  lf/.style={rectangle, rounded corners=2pt, draw=greenhl, line width=0.9pt,
             fill=greenhl!14, text=greenhl!50!black, font=\scriptsize\bfseries,
             minimum width=0.62cm, minimum height=0.52cm},
  sk/.style={rectangle, rounded corners=2pt, draw=neutral!45, line width=0.7pt,
             fill=gray!5, text=neutral!70, font=\scriptsize,
             minimum width=0.62cm, minimum height=0.52cm},
  ed/.style={neutral!70, line width=0.8pt},
  dead/.style={neutral!35, line width=0.8pt, dash pattern=on 2pt off 2pt},
  lab/.style={font=\tiny, text=neutral},
  cut/.style={accent, line width=1.2pt}]

\node[mx] (r) at (0,3.1) {5};
\node[mn] (m1) at (-2.6,1.9) {5};
\node[mn] (m2) at (2.6,1.9)  {2};
\node[mx] (a) at (-3.9,0.7) {5};
\node[mx] (b) at (-1.3,0.7) {6};
\node[mx] (c) at (1.3,0.7)  {2};
\node[sk] (d) at (3.9,0.7)  {?};
\node[lf] (l1) at (-4.55,-0.5) {3};
\node[lf] (l2) at (-3.25,-0.5) {5};
\node[lf] (l3) at (-1.95,-0.5) {6};
\node[sk] (l4) at (-0.65,-0.5) {9};
\node[lf] (l5) at (0.65,-0.5)  {1};
\node[lf] (l6) at (1.95,-0.5)  {2};
\node[sk] (l7) at (3.25,-0.5)  {0};
\node[sk] (l8) at (4.55,-0.5)  {-1};

\foreach \a/\b in {r/m1,r/m2,m1/a,m1/b,m2/c,
                   a/l1,a/l2,b/l3,c/l5,c/l6}{\draw[ed] (\a) -- (\b);}
\foreach \a/\b in {m2/d,b/l4,d/l7,d/l8}{\draw[dead] (\a) -- (\b);}

% cutoff 1: at the MAX node b, after leaf 6, alpha=6 >= beta=5
\draw[cut] (-1.10,0.28) -- (-0.60,-0.12);
\draw[cut] (-0.60,0.28) -- (-1.10,-0.12);
\node[font=\tiny\bfseries, text=accent, anchor=west] at (-0.48,0.10)
  {$\alpha=6 \ge \beta=5$};

% cutoff 2: at the MIN node m2, after child c, beta=2 <= alpha=5
\draw[cut] (3.05,1.45) -- (3.55,1.05);
\draw[cut] (3.55,1.45) -- (3.05,1.05);
\node[font=\tiny\bfseries, text=accent, anchor=west] at (3.65,1.25)
  {$\beta=2 \le \alpha=5$};

\node[lab, anchor=east] at (-5.1,3.1)  {\textbf{MAX} (us)};
\node[lab, anchor=east] at (-5.1,1.9)  {\textbf{MIN} (them)};
\node[lab, anchor=east] at (-5.1,0.7)  {\textbf{MAX} (us)};
\node[lab, anchor=east] at (-5.1,-0.5) {leaves};

\node[rounded corners=3pt, fill=greenhl!8, draw=greenhl!50, line width=0.9pt,
      text=greenhl!45!black, font=\scriptsize, text width=13.4cm,
      align=flush left, inner sep=6pt] at (0,-1.9)
  {\textbf{Five leaves evaluated instead of eight, and the answer is still 5.}
   The first cutoff: our node has already found $6$, and their node above has
   already secured $5$ --- so they will never come this way, and leaf~$9$ is
   irrelevant. The second is larger: their right-hand node has found $2$, and we
   have already secured $5$ at the root, so we will never come this way and the
   entire right subtree below it --- three more leaves --- is skipped. Note that
   the dashed node's value is printed as ``?'': alpha--beta never learns it, and
   does not need to.};
\end{tikzpicture}}%
{Alpha--beta on the tree of \dsref{game-tree}. Green leaves are evaluated, grey are
never looked at, and the crosses mark the two cutoffs.}{alphabeta}

\section{Move Ordering Is Everything}

Alpha--beta prunes when it has already found something good. So the order in which
moves are tried decides how much is pruned: try the best move first and the bounds
tighten immediately; try it last and there is nothing to prune against.

With \emph{perfect} ordering --- best move first at every node --- alpha--beta
evaluates
\[
  b^{\lceil d/2 \rceil} + b^{\lfloor d/2 \rfloor} - 1
\]
leaves instead of $b^{d}$. That is roughly $b^{d/2}$: the same time buys
\emph{twice the depth}, which for a game-playing program is the whole ball game.

The chapter's lab measures all of this on a five-ply, three-branch negotiation ---
$243$ leaves --- and the perfect-ordering bound is reached exactly:
$3^{3} + 3^{2} - 1 = 27 + 9 - 1 = 35$.

\smallskip
\noindent\begin{longtable}{@{}L{5.6cm}L{2.4cm}L{2.0cm}L{4.6cm}@{}}
\toprule
\rowcolor{primary!15}
\textbf{Algorithm} & \textbf{Leaves seen} & \textbf{Saving} & \textbf{Note} \\
\midrule
\endhead
Minimax, no pruning & 243 & --- & Every leaf, always \\
Alpha--beta, worst ordering & 217 & 1.1$\times$ & Pruning still helps a little,
even at its worst \\
Alpha--beta, natural ordering & 75 & 3.2$\times$ & The realistic case \\
Alpha--beta, perfect ordering & \textbf{35} & \textbf{6.9$\times$} & Exactly
$b^{\lceil d/2\rceil} + b^{\lfloor d/2\rfloor} - 1$ \\
\bottomrule
\caption{Leaf evaluations on a five-ply, three-branch negotiation. All four return
the same minimax value of $14$ --- pruning is exact, not approximate.}
\label{tab:alphabeta}
\end{longtable}

\begin{conceptbox}
Perfect ordering requires knowing the best move, which is what the search is for.
So the practical question is how good an ordering you can buy cheaply. Real engines
use a shallow search first and order the deep search by its results, plus
remembered ``killer'' moves that caused cutoffs elsewhere.

But cheaply is doing real work in that sentence. The lab measures a shallow-probe
ordering on this tree and it costs \emph{506} leaf evaluations against natural
ordering's $75$ --- nearly seven times worse, because the probes themselves are
evaluations and the tree is not deep enough to amortise them. Ordering pays in
chess, at twenty ply with a transposition table caching the probes. At five ply it
does not. Measure before you believe.
\end{conceptbox}

\section{Cutting Off the Search}

Real games do not fit. A negotiation with ten rounds and five options each is
$5^{10} \approx 10$ million leaves, and chess is unthinkable. So replace the
terminal test with a \term{cutoff test} --- stop at a fixed depth --- and replace
the utility with an \term{evaluation function} that estimates the position's worth.

The evaluation function is where domain knowledge enters, and the search is no
better than it. For the negotiation: how far the current offer is from our target,
how much slack remains in the schedule, whether the remaining rounds favour us.

\term{The horizon effect} is the price. A cutoff at depth $d$ cannot see anything at
depth $d+1$, so the search happily chooses a line that pushes an unavoidable loss
just past the horizon --- conceding a small point now to delay a large concession
into invisibility. The standard mitigations are to search \emph{quiescent}
positions deeper (do not stop in the middle of an exchange) and to extend the
search along forcing lines.

\begin{worked}[Tracing alpha--beta through the tree of \dsref{game-tree}]
Left to right, carrying $(\alpha, \beta)$.

\smallskip
\textbf{Root} (MAX), $\alpha = -\infty$, $\beta = +\infty$. Descend left.

\smallskip
\textbf{Left MIN node}, $(-\infty, +\infty)$. Descend to its left child.

\smallskip
\textbf{MAX node }$a$, $(-\infty, +\infty)$: evaluate leaf $3$, so $v=3$,
$\alpha = 3$; evaluate leaf $5$, so $v=5$, $\alpha=5$. Return $5$.
\emph{(2 leaves)}

\smallskip
Back at the left MIN node: $v = 5$, so $\beta = 5$. Descend to its right child
with $(-\infty, 5)$.

\smallskip
\textbf{MAX node }$b$, $(-\infty, 5)$: evaluate leaf $6$, so $v = 6$ and
$\alpha = 6$. Now $\alpha \ge \beta$ ($6 \ge 5$) --- \textbf{cutoff}. Leaf $9$ is
never evaluated. Return $6$. \emph{(1 leaf)}

\smallskip
Why it is safe: node $b$ is a MAX node that can already secure $6$. The MIN node
above it will never accept $6$ when it already has $5$ available. Whatever leaf $9$
holds, $b$ can only get \emph{better} than $6$, so MIN's choice is settled.

\smallskip
Left MIN node returns $\min(5, 6) = 5$. Root: $v = 5$, $\alpha = 5$.

\smallskip
\textbf{Right MIN node}, $(5, +\infty)$. Descend left.

\smallskip
\textbf{MAX node }$c$, $(5, +\infty)$: leaf $1$ gives $v=1$; leaf $2$ gives $v=2$.
Return $2$. \emph{(2 leaves)}

\smallskip
Right MIN node: $v = 2$, so $\beta = 2$. Now $\beta \le \alpha$ ($2 \le 5$) ---
\textbf{cutoff}. Node $d$ and both its leaves are skipped entirely. Return $2$.

\smallskip
Why it is safe: this MIN node can already hold us to $2$, and we have already
secured $5$ at the root. We will never play into this branch, so its exact value is
irrelevant --- which is why \dsref{alphabeta} prints it as ``?''.

\smallskip
\textbf{Root} returns $\max(5,2) = 5$. \textbf{Leaves evaluated: $3, 5, 6, 1, 2$
--- five of eight}, and the value is identical to minimax's.

\smallskip
\textbf{Now reverse the leaves} to $[-1, 0, 2, 1, 9, 6, 5, 3]$ and trace again. The
minimax value is still $5$, but the cutoffs fall differently --- that is move
ordering, and it is the entire content of Table~\ref{tab:alphabeta}.
\end{worked}

\begin{pitfall}
\begin{tight}
  \item \textbf{Thinking alpha--beta is an approximation.} It returns exactly the
        minimax value. What it approximates is nothing; what it saves is leaves.
  \item \textbf{Modelling a negotiator as a stochastic process.} If their choice
        depends on yours, an average-case model is not merely inaccurate, it is
        exploitable.
  \item \textbf{Quoting a cooperative-model number as a guarantee.} The minimax
        value is what you can secure \emph{whatever} they do. Any friendlier
        assumption gives a bigger number and no floor.
  \item \textbf{Assuming move ordering always pays.} Measured here: a shallow-probe
        ordering cost $506$ leaf evaluations against natural ordering's $75$. It
        pays at depth twenty, not at depth five.
  \item \textbf{Stopping the search in a volatile position.} The horizon effect
        makes the program push bad news just out of sight. Search quiescent
        positions deeper.
  \item \textbf{Forgetting that the evaluation function is the ceiling.} Beyond a
        point, deeper search of a bad evaluation function just finds worse moves
        more confidently.
\end{tight}
\end{pitfall}

%---------------------------------------------------------------------
\begin{lab}[Minimax and Alpha--Beta on a Slot Negotiation]
The five-ply negotiation of Table~\ref{tab:alphabeta}, with four move orderings,
counting leaf evaluations --- the operation that costs.
\begin{lstlisting}[language=Python]
"""Chapter 8 lab -- minimax and alpha-beta on a release negotiation.

Five plies, three options each, 243 leaves. All four configurations
return the same minimax value; they differ only in how many leaves
they have to look at. The perfect-ordering row hits the theoretical
bound b^ceil(d/2) + b^floor(d/2) - 1 exactly.
"""

import random

DEPTH = 5          # plies of negotiation
BRANCH = 3         # hold, concede, swap
OPTIONS = ("hold", "concede", "swap")


def build(seed=7):
    """Leaf utilities to us, in days of delay avoided."""
    rng = random.Random(seed)
    leaves = {}

    def rec(path):
        if len(path) == DEPTH:
            leaves[path] = rng.randrange(-20, 21)
            return
        for m in range(BRANCH):
            rec(path + (m,))

    rec(())
    return leaves


LEAVES = build()


class Counter:
    def __init__(self):
        self.evals = 0          # leaf evaluations: the expensive op
        self.nodes = 0


def minimax(path, maximising, c):
    c.nodes += 1
    if len(path) == DEPTH:
        c.evals += 1
        return LEAVES[path]
    vals = [minimax(path + (m,), not maximising, c)
            for m in range(BRANCH)]
    return max(vals) if maximising else min(vals)


def true_value(path, maximising):
    """An oracle, used only to build the best/worst orderings."""
    return minimax(path, maximising, Counter())


def ordered(path, maximising, how):
    ms = list(range(BRANCH))
    if how == "natural":
        return ms
    # 'best' puts the strongest move first, 'worst' puts it last
    rev = maximising if how == "best" else not maximising
    return sorted(ms, key=lambda m: true_value(path + (m,), not maximising),
                  reverse=rev)


def alphabeta(path, maximising, alpha, beta, c, how="natural"):
    c.nodes += 1
    if len(path) == DEPTH:
        c.evals += 1
        return LEAVES[path]
    if maximising:
        v = -10 ** 9
        for m in ordered(path, maximising, how):
            v = max(v, alphabeta(path + (m,), False, alpha, beta, c, how))
            alpha = max(alpha, v)
            if alpha >= beta:
                break                      # beta cutoff
        return v
    v = 10 ** 9
    for m in ordered(path, maximising, how):
        v = min(v, alphabeta(path + (m,), True, alpha, beta, c, how))
        beta = min(beta, v)
        if beta <= alpha:
            break                          # alpha cutoff
    return v


def shallow(path, maximising, limit, c):
    """A depth-limited probe, used to ORDER a deeper search.

    Every leaf it touches is charged, because in a real engine the
    evaluation function is what costs.
    """
    if len(path) == DEPTH:
        c.evals += 1
        return LEAVES[path]
    if limit == 0:
        p = path
        while len(p) < DEPTH:
            p = p + (0,)
        c.evals += 1
        return LEAVES[p]
    vals = [shallow(path + (m,), not maximising, limit - 1, c)
            for m in range(BRANCH)]
    return max(vals) if maximising else min(vals)


def ab_probe_ordered(path, maximising, alpha, beta, c, probe=1):
    """Alpha-beta ordering its moves by a shallow probe."""
    c.nodes += 1
    if len(path) == DEPTH:
        c.evals += 1
        return LEAVES[path]
    ms = sorted(range(BRANCH),
                key=lambda m: shallow(path + (m,), not maximising,
                                      probe, c),
                reverse=maximising)
    if maximising:
        v = -10 ** 9
        for m in ms:
            v = max(v, ab_probe_ordered(path + (m,), False, alpha, beta,
                                        c, probe))
            alpha = max(alpha, v)
            if alpha >= beta:
                break
        return v
    v = 10 ** 9
    for m in ms:
        v = min(v, ab_probe_ordered(path + (m,), True, alpha, beta,
                                    c, probe))
        beta = min(beta, v)
        if beta <= alpha:
            break
    return v


if __name__ == "__main__":
    print("depth %d, branching %d, leaves %d"
          % (DEPTH, BRANCH, BRANCH ** DEPTH))

    c0 = Counter()
    value = minimax((), True, c0)
    print("minimax value: %d" % value)
    print()
    print("algorithm                      leaf evals   nodes  value")
    print("-" * 60)
    print("%-30s %10d %7d %6d"
          % ("minimax, no pruning", c0.evals, c0.nodes, value))

    results = {"minimax": c0.evals}
    for how, label in (("worst", "alpha-beta, worst ordering  "),
                       ("natural", "alpha-beta, natural ordering"),
                       ("best", "alpha-beta, perfect ordering")):
        c = Counter()
        v = alphabeta((), True, -10 ** 9, 10 ** 9, c, how)
        assert v == value, how          # pruning is EXACT
        results[how] = c.evals
        print("%-30s %10d %7d %6d" % (label, c.evals, c.nodes, v))

    # the theoretical bound, reached exactly by perfect ordering
    bound = BRANCH ** ((DEPTH + 1) // 2) + BRANCH ** (DEPTH // 2) - 1
    print()
    print("b^ceil(d/2) + b^floor(d/2) - 1 = %d^%d + %d^%d - 1 = %d"
          % (BRANCH, (DEPTH + 1) // 2, BRANCH, DEPTH // 2, bound))
    assert results["best"] == bound
    print("perfect ordering evaluated %d leaves: the bound, exactly."
          % results["best"])

    # pruning never loses, and ordering decides how much it wins
    assert results["worst"] < results["minimax"]
    assert results["best"] < results["natural"] < results["worst"]
    print()
    print("saving against minimax:  worst %.1fx   natural %.1fx   "
          "perfect %.1fx"
          % (results["minimax"] / results["worst"],
             results["minimax"] / results["natural"],
             results["minimax"] / results["best"]))

    # does a shallow probe pay for the ordering it buys? Not here.
    c = Counter()
    v = ab_probe_ordered((), True, -10 ** 9, 10 ** 9, c, probe=1)
    assert v == value
    print()
    print("ordering by a depth-1 probe: %d leaf evals against natural"
          % c.evals)
    print("ordering's %d. The probes cost more than they save at this"
          % results["natural"])
    print("depth. Ordering pays in chess, at twenty ply with a cache.")
    assert c.evals > results["natural"]

    # how much does the instance matter?
    totals = []
    for seed in range(20):
        globals()["LEAVES"] = build(seed)
        c = Counter()
        alphabeta((), True, -10 ** 9, 10 ** 9, c, "natural")
        totals.append(c.evals)
    print()
    print("across 20 different negotiations, natural ordering:")
    print("   leaf evals  min %d   mean %.1f   max %d"
          % (min(totals), sum(totals) / len(totals), max(totals)))
\end{lstlisting}
Then try the experiment from the worked example: set \texttt{DEPTH = 6} and watch
both the minimax count and the perfect-ordering bound move --- $729$ leaves against
$3^{3}+3^{3}-1 = 53$. The gap is what buys the extra depth.
\end{lab}

%---------------------------------------------------------------------
\begin{checkpoint}
\begin{tightnum}
  \item Alpha--beta evaluated five leaves where minimax evaluated eight, and both
        returned $5$. Explain in one sentence why that is not a coincidence and not
        an approximation.
  \item In \dsref{alphabeta} the right-hand MIN node's right child is printed as
        ``?''. State the two facts that together make its value irrelevant.
  \item A colleague adds shallow-probe move ordering to your alpha--beta search and
        reports that it is slower. Give the most likely explanation, and say what
        would have to change for the ordering to pay.
\end{tightnum}
\end{checkpoint}

%---------------------------------------------------------------------
\begin{chaptersummary}
\begin{tight}
  \item A \term{game tree} alternates \term{MAX} levels (ours) and \term{MIN}
        levels (theirs). The \term{minimax value} is the position's worth under
        optimal play by both sides, and it is a \emph{guarantee}: at least this
        much, whatever they do.
  \item Minimax assumes a zero-sum optimal opponent. That is usually false and
        still the right default, because it is the only assumption that yields a
        floor.
  \item \term{Alpha--beta pruning} carries the best values each player can already
        secure and stops examining a branch when $\alpha \ge \beta$. It returns
        \emph{exactly} the minimax value.
  \item \term{Move ordering} decides how much is pruned. With perfect ordering
        alpha--beta sees $b^{\lceil d/2\rceil} + b^{\lfloor d/2\rfloor} - 1$
        leaves --- measured here as exactly $35$ against minimax's $243$ ---
        which is roughly $b^{d/2}$, so the same time buys twice the depth.
  \item Buying a good ordering is not free. A shallow-probe ordering cost $506$
        leaf evaluations here against natural ordering's $75$: it pays at great
        depth with a cache, and not at five ply.
  \item Real games need a \term{cutoff test} and an \term{evaluation function},
        which is where domain knowledge enters and where the search's ceiling is
        set. The \term{horizon effect} is the cost: the program pushes
        unavoidable bad news just past the cutoff.
\end{tight}
\end{chaptersummary}

\begin{reviewq}
  \item \textbf{(MCQ)} Alpha--beta pruning returns: (a)~an approximation of the
        minimax value (b)~exactly the minimax value (c)~the greedy value
        (d)~a lower bound only.
  \item \textbf{(MCQ)} A cutoff occurs when: (a)~$\alpha \ge \beta$
        (b)~$\alpha < \beta$ (c)~the depth limit is reached (d)~a leaf is
        negative.
  \item \textbf{(MCQ)} Under perfect move ordering alpha--beta examines about:
        (a)~$b^{d}$ (b)~$b^{d/2}$ (c)~$d^{b}$ (d)~$bd$ leaves.
  \item \textbf{(MCQ)} The horizon effect is: (a)~running out of memory
        (b)~pushing an unavoidable loss beyond the depth limit (c)~pruning too
        aggressively (d)~a bad evaluation function.
  \item \textbf{(Short)} Explain why the minimax value is a guarantee, and what you
        give up by modelling the other party as cooperative instead.
  \item \textbf{(Short)} Why does move ordering affect alpha--beta but not minimax?
  \item \textbf{(Short)} Using the measured numbers in this chapter, explain the
        circumstances in which investing in move ordering pays and the
        circumstances in which it does not.
  \item \textbf{(Applied)} Take the tree of \dsref{game-tree} with its leaves
        reversed to $[-1, 0, 2, 1, 9, 6, 5, 3]$. Compute the minimax value, then
        trace alpha--beta left to right and list which leaves are evaluated and
        where each cutoff occurs. Compare the count with the five of the worked
        example and say what that demonstrates.
\end{reviewq}
